The F.3 resolution rule establishes when block-group resolution is required for
congressional redistricting: when $k/n > 0.05$ (where $k$ is the number of congressional
seats and $n$ is the number of census tracts in the state), tract-level redistricting
produces insufficient internal partitioning choices for the recursive bisection algorithm
to achieve population balance without excessive within-tract splits. Applying this rule
to Census Bureau medium-series 2030 population projections and projected Huntington-Hill
reapportionment outcomes, we find that 43 states will require block-group resolution
in 2030, up from 39 states in 2020. The four newly upgraded states are Montana ($k$: $1\to2$),
Idaho ($k$: $2\to3$), Colorado ($k$: $8\to9$), and North Carolina ($k$: $14\to15$).

Block-group adjacency graphs are approximately 5$\times$ larger than tract adjacency
graphs ($\sim$1.2 million vertices vs.\ $\sim$240{,}000 vertices nationally). The additional
scale affects three operational parameters: construction time (estimated 20 CPU-hours
vs.\ 4 CPU-hours for tracts), memory footprint per state (estimated peak 14 GB for Texas
vs.\ 2.8 GB at tract resolution), and full-sweep redistricting runtime (estimated 90 minutes
at block-group resolution vs.\ 90 minutes at tract resolution, since the bisection time
scales with adjacency complexity, not raw vertex count).

Pre-building the 43-state block-group adjacency by January 2031 reduces the effective
2031 redistricting window from 6 months (if adjacency construction occurs during
redistricting) to 2 months (loading pre-built graphs and running bisection). This
4-month reduction is the operational motivation for the Q.3 infrastructure investment.
The \textsc{bisect} CLI command for block-group adjacency construction is:

\begin{verbatim}
bisect fetch --year 2030 --resolution block-group --workers 8
\end{verbatim}

We document the state-by-state resolution classification, construction resource
estimates, and validation protocols for the 43-state block-group adjacency build.
